\documentclass[12font]{article}
\title{\textbf{ Time/Temperature Equivalence }}
\author{Chen-Chih Wang  }
\date{\textit{Department of physics, National Tsing Hua University, Hsinchu 30013, Taiwan} \\ (Dated: \today)}
\usepackage[margin=2cm]{geometry}
%\usepackage[fleqn]{amsmath}
%\usepackage{CJK,amsmath}
\usepackage{amsfonts,amssymb}
\usepackage{fancyhdr}
%\pagestyle{fancy}
\usepackage{textcomp}
\usepackage{graphicx}
\usepackage{float}
\usepackage{color}
\usepackage{tikz}
\usepackage{multicol}
\usepackage{threeparttable}
\usepackage{amsthm}
\usepackage{mathtools}
\usepackage{hyperref}
\urlstyle{same}
\usepackage{caption}
\usepackage{subcaption}

\newcommand{\bra}[1]{\left\langle #1 \right|}
\newcommand{\ket}[1]{\left| #1 \right\rangle}
\newcommand{\anglelr}[1]{\left\langle #1 \right\rangle}
\newcommand{\braket}[2]{\left\langle #1 \right|\left. #2 \right\rangle}
\newcommand{\dif}[1]{\frac{\mathrm{d}}{\mathrm{d}#1}}
\newcommand{\diff}[2]{\frac{\mathrm{d}#1}{\mathrm{d}#2}}
\newcommand{\gd}{\boldsymbol{\nabla}}
\newcommand{\dive}{\boldsymbol{\nabla}\cdot}
\newcommand{\curl}{\boldsymbol{\nabla}\times}
\newcommand{\lr}[1]{\left(  #1  \right)}
\newcommand{\lrm}[1]{\left[  #1  \right]}
\newcommand{\lrg}[1]{\left\{  #1  \right\}}
\newcommand{\abs}[1]{\left|  #1  \right|}
\newcommand{\de}{\partial}
\newcommand{\pdif}[1]{\frac{\partial}{\partial#1}}
\newcommand{\pdiff}[2]{\frac{\partial#1}{\partial#2}}
\newcommand{\mrm}[1]{\mathrm{#1}}
\newcommand{\bo}[1]{\mathbf{#1}}
\newcommand{\bog}[1]{\boldsymbol{#1}}
\newcommand{\vint}[3]{\int_{#1}#2\cdot\mathrm{d}\mathbf{#3}}
\newcommand{\voint}[3]{\oint_{#1}#2\cdot\mathrm{d}\mathbf{#3}}
\newcommand{\Vint}{\int\mrm{d}^3\bo{r}'}
\newcommand{\rr}{\mathbf{r} - \mathbf{r}'}
\newcommand{\arr}{\left|\mathbf{r} - \mathbf{r}'\right|}

\begin{document}
\setlength{\baselineskip}{20pt}

\maketitle
\begin{figure}[H] %abstract
  \centering
  \begin{minipage}[t]{0.8\textwidth}
  \small{
  \centering
     \textbf{Abstract} 
     
     Feynman path integral formalism is useful for both quantum field theory and the statistical physics. The partition function can be reformulated as a path integral, and the inverse of the temperature plays the role as the time in the quantum field theory (which is so called the imaginary time). It has been proposed that there might be a correspondence between the statistical mechanics and the quantum mechanics, based on the similarity of the two important factors, the Boltzmann factor $e^{-\hat{H}/T}$ in the statistical mechanics and the time evolution operator $e^{-i\hat{H}t}$ in the quantum mechanics. Here I try to find more evidence to prove that there indeed exist a time/temparature correspondence, not only mathematically, but trying to find the physical meaning of this similarity.
  }
  \end{minipage}
\end{figure}

\section{Introduction}
Feynmann path integral formalism is one of the quantization methods. It is convenient for taking the perturbation treatment using the Feynman diagram technique. In addition, it also provides a quantum mechanics point of view to understand the least action principle. The time evolution operator in the Feynman path integral is
\[\hat{U}(t_f,t_i) = e^{-i\hat{H}t} = \int \mathcal{D}[\psi,\bar{\psi}]e^{i\mathcal{S}[\psi,\bar{\psi}]} \]
On the other hand, the partition function can also be written in the path integral form
\[\mathcal{Z} = \mrm{tr}\lr{e^{-\hat{H}/T}} = \int \mathcal{D}[\psi,\bar{\psi}]e^{-\mathcal{S}[\psi,\bar{\psi}]}\]
where in the statistical mechanics
\[
\mathcal{S}[\psi,\bar{\psi}] = \int_{0}^{1/T}\lr{}
\]

\section{Similarity between the }


%\begin{thebibliography}{99}

%    \bibitem{}{}

%\end{thebibliography}

\end{document}

